\begin{helper}[Nonroot interchangeable components are cycles]
Let \(A\) be an initial segment of a tree sequence \(B\), and let
\(\varphi\) be a proper partial edge coloring.  Suppose \(\alpha\ne\beta\)
implicitly through the assertion that the component through a designated
vertex \(s\in V(A)\) is a two-color path component.  Suppose that this
root component meets \(A\) only at \(s\), that \(\alpha,\beta\) are
\(B\)-interchangeable, and that every component rooted at a vertex of
\(A-s\) contains an exact edge-node.  Then every such nonroot component
is a full two-color cycle and avoids \(s\).
\end{helper}
\begin{proof}
Fix \(v\in V(A)-\{s\}\), and let \(C_v\) be its full
\((\alpha,\beta)\)-component.  The initial-segment relation
\noderef{TreeSequenceInitialSegment} puts both \(s\) and \(v\) in
\(V(B)\).  If \(C_v\) were a \noderef{TwoColorPathComponent}, then
\noderef{TreeInterchangeable}, applied to \(C_v\) and the designated
root path component, would identify their full supports because both
meet \(B\).  The vertex-node of \(v\) would then lie in the root
component, contradicting the hypothesis that its only vertex in \(A\)
is \(s\).  Hence \(C_v\) is not a path component.

By hypothesis, \(C_v\) contains an exact edge-node, so it is not the
isolated-component alternative.  The root path-component hypothesis
also includes \(\alpha\ne\beta\).  Unfolding
\noderef{TwoColorPathComponent}, failure of the path alternative now
says that no node of \(C_v\) has a subsingleton neighbor set.  On the
other hand, \noderef{TwoColorIncidenceDegreeAtMostTwo}, applied to the
proper coloring \(\varphi\), gives degree at most two at every node.
Thus every node of \(C_v\) has exactly two neighbors.

Take the simple spanning path supplied by
\noderef{TwoColorComponentSpanningPath}; its support is exactly \(C_v\).
Each endpoint has one path neighbor and therefore has one further
neighbor in \(C_v\).  That further neighbor cannot be an internal path
node, since an internal node already has its two path neighbors and
would acquire degree at least three.  It is therefore the other endpoint.
Adding this last adjacency closes the spanning path into a simple closed
walk, hence a \noderef{TwoColorCycle}, without changing its support.
The support equivalence is exactly
\noderef{TwoColorComponentCycle} for \(C_v\).

Finally, if the vertex-node of \(s\) lay in \(C_v\), connected-component
maximality in \noderef{TwoColorComponent} would identify \(C_v\) with
the root component.  That would again put \(v\) in the root component,
contrary to its asserted intersection with \(A\).  Thus \(C_v\) avoids
\(s\), as required.
\end{proof}
